\section{The fractional optimum is bounded by the tree}\label{sec:lp}

This section proves that the cost of any invariant perfect matching on the Poisson
hypertree is an asymptotic upper bound for the fractional deficit of the market. The tree
and invariance are those of Section~\ref{sec:tree}; balls, the truncated tree $\tree_K$ and
the bound $\eps_1(q,K,r)$ are those of Section~\ref{sec:localconv}; $C_0$ is the constant of
\eqref{eq:variance-C0}.

\subsection{Statement}\label{sec:lp-statement}

We recall the setting of Section~\ref{sec:tree}. The Poisson hypertree $(\tree,o,\xi)$ of
Definition~\ref{def:tree-poisson} is a random rooted hypertree with costs; its root sector
is uniform; almost surely every agent lies in infinitely many teams, with pairwise distinct
costs of which finitely many lie below any level. Rooted marked hypertrees are coded by
ranks, and the $\sigma$-field on them is generated by the truncated balls $B_r(\tree_K,o)$,
$r,K\in\N$ (Remark~\ref{rem:tree-balls}; $\tree_K$ is the truncation written $T_K$ there,
applied to $\tree$, and the balls are those of Definition~\ref{def:localconv-ball}). An
\emph{invariant perfect matching} on $\tree$ is a family $m=(m_e)$ of marks in $\{0,1\}$ on
the teams of $\tree$, possibly defined with the help of additional randomness, such that
almost surely every agent lies in exactly one team $e$ with $m_e=1$, and such that the
decorated tree $(\tree,o,\xi,m)$ is invariant in the sense of
Definition~\ref{def:tree-invariant}. Two agents are adjacent, $a\sim b$, when they lie in a
common team, which is then unique. Invariance is the identity \eqref{eq:tree-MT}: for every
measurable $f\ge0$ on doubly rooted decorated hypertrees,
\begin{equation}\label{eq:lp-mt}
\E\sum_{b\sim o}f(o,b)=\E\sum_{b\sim o}f(b,o),
\end{equation}
where $f(a,b)$ is the value of $f$ at the decorated tree with first root $a$ and second
root $b$. By linearity \eqref{eq:lp-mt} also holds for signed $f$ with
$\E\sum_{b\sim o}(|f(o,b)|+|f(b,o)|)<\infty$. The \emph{cost} \eqref{eq:tree-cost} of $m$ is
\[
c(m)=\E\sum_{e\ni o}\cost_e\,m_e\ \in[0,\infty],
\]
the expected cost of the team of the root.

\begin{theorem}\label{thm:lp}
Let $m$ be an invariant perfect matching on $\tree$. Then
\[
\limsup_{q\to\infty}\frac{\E\PhiLP}{q^{1/3}}\ \le\ c(m).
\]
More precisely, there are numbers $c_r\le c(m)$ and $\eps_r\in[0,1]$, $r\ge1$, depending only
on the law of $(\tree,o,\xi,m)$, with $\eps_r\to0$ as $r\to\infty$, such that for all
integers $r\ge1$ and $q\ge r^{3/2}$ and every $\eta\in(0,1)$,
\begin{equation}\label{eq:lp-explicit}
\frac{\E\PhiLP}{q^{1/3}}\ \le\ \frac{c_r+r\,\eps_1(q,r,r+1)}{1-\eta}
+\frac{C_0}{\eta}\bigl(\eps_r+\eps_1(q,r,r+1)\bigr)^{1/3},
\end{equation}
with $\eps_1$ as in Theorem~\ref{thm:localconv}.
\end{theorem}

If $c(m)=\infty$ there is nothing to prove, and we assume $c(m)<\infty$. The proof has five
steps. The conditional expectation of $m$ given the tree is a fractional perfect matching
that all three agents of a team compute alike (Lemma~\ref{lem:lp-conditional}). Conditioning
further on a ball of finite radius, and taking the minimum over the three agents of a team,
gives a fractional packing of finite range whose cost is at most $c(m)$ and whose slack is
small (Lemmas~\ref{lem:lp-tree} and~\ref{lem:lp-condrule}). The same rule is applied in the
market, where its mean cost and mean slack are those of the tree up to the total variation
distance of Theorem~\ref{thm:localconv} (Lemma~\ref{lem:lp-market}). The pair deficits are
split into two independent uniform families (Lemma~\ref{lem:lp-split}): the packing is built
on the first, and its slack is filled by optimal assignments of the second on random sets of
agents, priced by $C_0$ (Lemma~\ref{lem:lp-patch}).

\subsection{The conditional matching}\label{sec:lp-conditional}

Let $\ell_1,\ell_2,\dots$ be the teams at the root $o$ in increasing order of cost. For an
agent $v$ and a team $e\ni v$ let $\operatorname{rk}_v(e)$ be the rank of $e$ among the teams
at $v$, so that $\operatorname{rk}_o(\ell_j)=j$. Fix measurable functions $h_j$, $j\ge1$, of
rooted marked hypertrees, with values in $[0,1]$, such that
\[
h_j(\tree,o,\xi)=\E\bigl[m_{\ell_j}\bigm|\tree,o,\xi\bigr]\qquad\text{almost surely}.
\]
For an agent $v$ and a team $e\ni v$ put
\[
x^v_e=h_{\operatorname{rk}_v(e)}(\tree,v,\xi),
\]
where $(\tree,v,\xi)$ is the tree re-rooted at $v$. So $x^v_e$ is the conditional
probability that $e$ is matched, computed from the point of view of $v$.

\begin{lemma}\label{lem:lp-conditional}
Almost surely $\sum_{e\ni o}x^o_e=1$, and $x^b_e=x^o_e$ for every team $e\ni o$ and every
agent $b\in e$. Moreover $\E\sum_{e\ni o}\cost_e\,x^o_e=c(m)$.
\end{lemma}

\begin{proof}
By monotone convergence for conditional expectations,
$\sum_jh_j=\E[\sum_{e\ni o}m_e\mid\tree,o,\xi]=1$ and
$\E\sum_j\cost_{\ell_j}h_j=\E\sum_j\cost_{\ell_j}m_{\ell_j}=c(m)$, because $\cost_{\ell_j}\ge0$
is a function of $(\tree,o,\xi)$.

For adjacent agents $a,b$ with common team $e$ put
\[
f(a,b)=\operatorname{sgn}\bigl(x^a_e-x^b_e\bigr)\,(m_e-x^b_e).
\]
This is a measurable function of the doubly rooted decorated tree, and
$|f(a,b)|\le m_e+x^b_e$. We have $\E\sum_{b\sim o}m_e=2$, because the matched team of the
root has two further agents, and $\E\sum_{b\sim o}x^o_e=2\,\E\sum_jh_j=2$. By
\eqref{eq:lp-mt} for the nonnegative function $(a,b)\mapsto x^b_e$,
$\E\sum_{b\sim o}x^b_e=\E\sum_{b\sim o}x^o_e=2$. So
$\E\sum_{b\sim o}(|f(o,b)|+|f(b,o)|)\le8$, and \eqref{eq:lp-mt} holds for $f$. We compute
both sides.

\emph{Right-hand side.} Grouping the neighbours of $o$ by their team,
\[
\E\sum_{b\sim o}f(b,o)=\E\sum_{j\ge1}\Psi_j\,(m_{\ell_j}-h_j),\qquad
\Psi_j=\sum_{b\in\ell_j\setminus\{o\}}\operatorname{sgn}\bigl(x^b_{\ell_j}-h_j\bigr).
\]
The variable $\Psi_j$ is a function of $(\tree,o,\xi)$, since the tree re-rooted at a
neighbour of the root is a measurable function of the tree (Section~\ref{sec:tree}); and $|\Psi_j|\le2$. Hence
$\E[\Psi_j(m_{\ell_j}-h_j)]=0$ for every $j$ by the definition of $h_j$. The sum over $j$
may be taken outside the expectation because $\E\sum_j|\Psi_j|(m_{\ell_j}+h_j)\le4$. So the
right-hand side is $0$.

\emph{Left-hand side.} For $b\in\ell_j\setminus\{o\}$ the variables $x^b_{\ell_j}$ and
$\operatorname{sgn}(h_j-x^b_{\ell_j})$ are functions of $(\tree,o,\xi)$. Conditioning on
$(\tree,o,\xi)$ replaces $m_{\ell_j}$ by $h_j$; this is justified by the integrability shown
above. Hence
\[
\E\sum_{b\sim o}f(o,b)=\E\sum_{j\ge1}\ \sum_{b\in\ell_j\setminus\{o\}}
\operatorname{sgn}\bigl(h_j-x^b_{\ell_j}\bigr)\bigl(h_j-x^b_{\ell_j}\bigr)
=\E\sum_{e\ni o}\ \sum_{b\in e\setminus\{o\}}\bigl|x^o_e-x^b_e\bigr| .
\]

The two sides are equal, so the last expectation is $0$.
\end{proof}

So $x_e:=x^o_e$, $e\ni o$, is a fractional perfect matching at the root that is a function
of the tree and its costs alone, has the cost of $m$, and is seen equally by the three
agents of each team. The matching $m$ itself may use additional randomness; from here on it
enters only through the functions $h_j$.

\subsection{Rules of finite range}\label{sec:lp-rules}

Fix an integer $r\ge1$ and a level $K>0$. We write $\mathcal G_r(K)$ for the space of balls
$\mathcal G_r$ of Section~\ref{sec:localconv-setting} when the level has to be shown.

\begin{definition}\label{def:lp-rule}
A \emph{rule of range $r$ at level $K$} is a sequence $\gamma=(\gamma_j)_{j\ge1}$ of
measurable functions $\gamma_j\colon\mathcal G_r(K)\to[0,1]$ such that, for every ball $B$:
$\sum_j\gamma_j(B)\le1$; $\gamma_j(B)=0$ if the root of $B$ lies in fewer than $j$ teams;
and $\gamma_j(B)=0$ for all $j$ unless $B$ is a hypertree in which the teams at every agent
have pairwise distinct costs.
\end{definition}

Let $G$ be a hypergraph of teams with costs in $[0,K]$ as in
Definition~\ref{def:localconv-ball}; below $G$ is $H_q(K)$ or $\tree_K$. A rule $\gamma$
assigns to every agent $v$ and every team $e\ni v$ of $G$ the number
\begin{equation}\label{eq:lp-y}
y^v_e=\gamma_{\operatorname{rk}_v(e)}\bigl(B_r(G,v)\bigr),
\end{equation}
where the rank is taken among the teams of $G$ at $v$; when two of these have equal costs
the rule vanishes and no rank is needed. The associated weights, slacks and local costs are
\begin{equation}\label{eq:lp-minpacking}
x^\gamma_e=\min_{v\in e}y^v_e,\qquad
\sigma_v=1-\sum_{e\ni v}x^\gamma_e,\qquad
\kappa_v=\sum_{e\ni v}\cost_e\,x^\gamma_e .
\end{equation}
Each agent proposes a weight for each of its teams from what it sees within distance $r$,
and a team receives the least of its three proposals.

\begin{lemma}\label{lem:lp-rule}
Let $\gamma$ be a rule of range $r$ at level $K$.
\begin{enumerate}
\item[\textup{(a)}] \emph{(Packing.)} $x^\gamma\ge0$ and
$\sum_{e\ni v}x^\gamma_e\le\sum_{e\ni v}y^v_e\le1$ for every agent $v$. Hence
$\sigma_v\in[0,1]$ and $\kappa_v\in[0,K]$.
\item[\textup{(b)}] \emph{(Range.)} There are measurable functions
$S^\gamma\colon\mathcal G_{r+1}(K)\to[0,1]$ and $C^\gamma\colon\mathcal G_{r+1}(K)\to[0,K]$,
depending only on $\gamma$, such that for every $G$ and every agent $v$ of $G$,
\[
\sigma_v=S^\gamma\bigl(B_{r+1}(G,v)\bigr),\qquad\kappa_v=C^\gamma\bigl(B_{r+1}(G,v)\bigr).
\]
\end{enumerate}
\end{lemma}

\begin{proof}
(a) $\sum_{e\ni v}y^v_e=\sum_j\gamma_j(B_r(G,v))\le1$, the sum being over the ranks of the
teams at $v$.

(b) Let $v$ be an agent and $e\ni v$ a team of $G$; it belongs to $B_{r+1}(G,v)$. Let
$b\in e$. A team of $B_r(G,b)$ contains an agent $a$ at distance at most $r-1$ from $b$,
hence at distance at most $r$ from $v$, so it is a team of $B_{r+1}(G,v)$. The same holds
for every team of a shortest chain from $b$ to $a$. Therefore $B_r(G,b)$ is the ball of
radius $r$ around $b$ inside the hypergraph $B_{r+1}(G,v)$, and it contains all teams of $G$
at $b$, so the rank of $e$ at $b$ can be read in it. Thus $x^\gamma_e$, for $e\ni v$, is a
function of $B_{r+1}(G,v)$. Define $S^\gamma(B')$ and $C^\gamma(B')$ for an arbitrary ball
$B'\in\mathcal G_{r+1}(K)$ with root $\rho$ by the formulas \eqref{eq:lp-y} and
\eqref{eq:lp-minpacking} for $v=\rho$, applied inside the hypergraph $B'$. The bounds follow
from (a).
\end{proof}

\subsection{The min-packing on the tree}\label{sec:lp-tree}

Apply a rule $\gamma$ of range $r$ at level $K$ to $G=\tree_K$. Truncation at $K$ removes
the most expensive teams at each agent, so it does not change the ranks of the remaining
teams. For a team $e$ of $\tree$ with $\cost_e>K$ we put $y^v_e=0$ for $v\in e$ and
$x^\gamma_e=0$; for a team with $\cost_e\le K$ at a neighbour $b$ of the root, $y^b_e$ is
given by \eqref{eq:lp-y} with $G=\tree_K$. We write
\[
\hat\gamma_j=y^o_{\ell_j}=\gamma_j\bigl(B_r(\tree_K,o)\bigr),
\]
which is $0$ when $\cost_{\ell_j}>K$.

\begin{lemma}\label{lem:lp-tree}
For every rule $\gamma$ of range $r$ at level $K$,
\begin{align*}
\E\,C^\gamma\bigl(B_{r+1}(\tree_K,o)\bigr)&\le\E\sum_{j\ge1}\cost_{\ell_j}\hat\gamma_j,\\
\E\,S^\gamma\bigl(B_{r+1}(\tree_K,o)\bigr)&\le3\,\E\sum_{j\ge1}\bigl(h_j-\hat\gamma_j\bigr)_+ .
\end{align*}
\end{lemma}

\begin{proof}
The first inequality holds because $x^\gamma_e\le y^o_e$ for $e\ni o$. For the second, let
$x_e=x^o_e$ be as in Lemma~\ref{lem:lp-conditional}. Almost surely $\sum_{e\ni o}x_e=1$ and
$x^v_e=x_e$ for all $e\ni o$ and $v\in e$, so
\[
\sigma_o=\sum_{e\ni o}\bigl(x_e-x^\gamma_e\bigr)
=\sum_{e\ni o}\max_{v\in e}\bigl(x^v_e-y^v_e\bigr)
\le\sum_{e\ni o}\ \sum_{v\in e}\bigl(x^v_e-y^v_e\bigr)_+ .
\]
The terms with $v=o$ sum to $\sum_j(h_j-\hat\gamma_j)_+$. For the others apply
\eqref{eq:lp-mt} to the nonnegative function $f(a,b)=(x^b_e-y^b_e)_+$, where $e$ is the
common team of $a$ and $b$; it is a measurable function of the doubly rooted tree. This
gives
\[
\E\sum_{b\sim o}\bigl(x^b_e-y^b_e\bigr)_+=\E\sum_{b\sim o}\bigl(x^o_e-y^o_e\bigr)_+
=2\,\E\sum_{j\ge1}\bigl(h_j-\hat\gamma_j\bigr)_+ ,
\]
since every team at the root contains two neighbours of the root.
\end{proof}

From now on the level is tied to the range: $K=K_r=r$. Let
\[
\mathcal F_r=\sigma\bigl(B_r(\tree_r,o)\bigr).
\]

\begin{lemma}\label{lem:lp-condrule}
For every $r\ge1$ there is a rule $g_r=(g_{r,j})_{j\ge1}$ of range $r$ at level $r$ such
that, almost surely, for every $j$,
\begin{equation}\label{eq:lp-condrule}
g_{r,j}\bigl(B_r(\tree_r,o)\bigr)=\E\bigl[m_{\ell_j}\ind\{\cost_{\ell_j}\le r\}\bigm|
\mathcal F_r\bigr]=\ind\{\cost_{\ell_j}\le r\}\,\E\bigl[h_j\bigm|\mathcal F_r\bigr].
\end{equation}
For these rules, with $c_r:=\E\,C^{g_r}(B_{r+1}(\tree_r,o))$ and
$\eps_r:=\E\,S^{g_r}(B_{r+1}(\tree_r,o))$,
\begin{enumerate}
\item[\textup{(a)}] $c_r\le\E\sum_{e\ni o}\cost_e\,m_e\ind\{\cost_e\le r\}\le c(m)$ for
every $r$;
\item[\textup{(b)}] $g_{r,j}(B_r(\tree_r,o))\to h_j$ almost surely as $r\to\infty$, for
every $j$, and $\eps_r\to0$.
\end{enumerate}
\end{lemma}

\begin{proof}
The event $\{\cost_{\ell_j}\le r\}$ says that the root of $B_r(\tree_r,o)$ lies in at least
$j$ teams, so it belongs to $\mathcal F_r$; this gives the second equality in
\eqref{eq:lp-condrule}, because $\mathcal F_r$ is contained in the $\sigma$-field of
$(\tree,o,\xi)$. Choose measurable functions $G_j\ge0$ on $\mathcal G_r(r)$ with
$G_j(B_r(\tree_r,o))=\E[m_{\ell_j}\ind\{\cost_{\ell_j}\le r\}\mid\mathcal F_r]$ almost surely.
Almost surely $\sum_jG_j\le1$, because $\sum_jm_{\ell_j}=1$; $G_j=0$ when the root lies in
fewer than $j$ teams; and the ball is a hypertree with pairwise distinct costs. Let $g_r$ be
equal to $(G_j)_j$ on the measurable set of balls where the three requirements of
Definition~\ref{def:lp-rule} hold for $(G_j)_j$, and $0$ elsewhere. The change is on a null
set of the law of $B_r(\tree_r,o)$.

(a) By Lemma~\ref{lem:lp-tree}, $c_r\le\E\sum_j\cost_{\ell_j}\,g_{r,j}(B_r(\tree_r,o))$.
The variable $\cost_{\ell_j}\ind\{\cost_{\ell_j}\le r\}$ is bounded and
$\mathcal F_r$-measurable, so by \eqref{eq:lp-condrule}
\[
\E\bigl[\cost_{\ell_j}\,g_{r,j}(B_r(\tree_r,o))\bigr]
=\E\bigl[\cost_{\ell_j}\ind\{\cost_{\ell_j}\le r\}\,m_{\ell_j}\bigr].
\]
Sum over $j$. The truncation costs nothing here: the rule ignores the teams above the level,
and what that loses is slack.

(b) The $\sigma$-fields $\mathcal F_r$ increase with $r$: deleting from
$B_{r+1}(\tree_{r+1},o)$ the teams of cost above $r$, keeping the component of the root and
then the teams that contain an agent at distance at most $r-1$ from the root, gives
$B_r(\tree_r,o)$. They generate the $\sigma$-field of $(\tree,o,\xi)$: by Remark~\ref{rem:tree-balls}
it is generated by the truncated balls $B_r(\tree_K,o)$ with integers $r,K$, and such a
ball is in the same way a function of $B_{r'}(\tree_{r'},o)$ with $r'=\max(r,K)$. By
L\'evy's upward theorem \cite{Kallenberg21},
$\E[h_j\mid\mathcal F_r]\to h_j$ almost surely, and
$\ind\{\cost_{\ell_j}\le r\}\to1$. This is the first claim. By Lemma~\ref{lem:lp-tree},
\[
\eps_r\le3\,\E\sum_{j\ge1}\bigl(h_j-g_{r,j}(B_r(\tree_r,o))\bigr)_+ .
\]
The $j$-th term of the series tends to $0$ almost surely and is at most $h_j$, and
$\sum_jh_j=1$. By dominated convergence for series the sum tends to $0$ almost surely; it is
at most $1$, so its expectation tends to $0$.
\end{proof}

The slack caused by the truncation at level $r$ is included in $\eps_r$; there is no
separate limit in $K$.

\subsection{Transfer to the market}\label{sec:lp-market}

\begin{lemma}\label{lem:lp-market}
Let $\gamma$ be a rule of range $r$ at level $K\le q^{2/3}$. Apply it to the near-tight
hypergraph $H=H_q(K)$ of the market, put $x^\gamma_e=0$ for teams that are not near-tight,
and let $V$ be a uniformly random agent independent of the market.
\begin{enumerate}
\item[\textup{(a)}] $x^\gamma$ is a function of the pair deficits; $x^\gamma\ge0$; every
agent has load $\sum_{e\ni v}x^\gamma_e\le1$; and $x^\gamma_e=0$ unless $\cost_e\le K$.
\item[\textup{(b)}] The total slack $\sum_v\sigma_v$ over the agents of a sector is the same
for the three sectors, namely $\Sigma:=q-\sum_ex^\gamma_e$.
\item[\textup{(c)}] The expected slack and the expected deficit of $x^\gamma$ are
\[
\E\,\Sigma=q\,\E\,S^\gamma\bigl(B_{r+1}(H,V)\bigr),\qquad
\E\sum_e\Defi(e)\,x^\gamma_e=q^{1/3}\,\E\,C^\gamma\bigl(B_{r+1}(H,V)\bigr).
\]
\item[\textup{(d)}] With the root sector of $\tree$ uniform,
\begin{align*}
\bigl|\E\,S^\gamma(B_{r+1}(H,V))-\E\,S^\gamma(B_{r+1}(\tree_K,o))\bigr|&\le\eps_1(q,K,r+1),\\
\bigl|\E\,C^\gamma(B_{r+1}(H,V))-\E\,C^\gamma(B_{r+1}(\tree_K,o))\bigr|&\le K\,\eps_1(q,K,r+1).
\end{align*}
\end{enumerate}
\end{lemma}

\begin{proof}
(a) is Lemma~\ref{lem:lp-rule}(a); the rule is defined on every ball, so nothing has to be
assumed about the realisation of the market.

(b) Every team has exactly one agent in each sector, so the total load of the agents of a
sector is $\sum_ex^\gamma_e$.

(c) Sums over $v$ are over all $3q$ agents. By (b) and Lemma~\ref{lem:lp-rule}(b),
\[
3\Sigma=\sum_v\sigma_v=\sum_vS^\gamma\bigl(B_{r+1}(H,v)\bigr),
\]
and the expectation of the last sum is $3q\,\E\,S^\gamma(B_{r+1}(H,V))$. In the same way,
since every team is counted at its three agents,
\[
\sum_vC^\gamma\bigl(B_{r+1}(H,v)\bigr)=\sum_v\kappa_v=3\sum_e\cost_e\,x^\gamma_e
=3q^{2/3}\sum_e\Defi(e)\,x^\gamma_e ,
\]
and the expectation of the first sum is $3q\,\E\,C^\gamma(B_{r+1}(H,V))$.

(d) For a measurable function $h$ with values in $[0,M]$ and probability measures $\mu$,
$\nu$, $|\int h\dd\mu-\int h\dd\nu|\le M\,d_{\mathrm{TV}}(\mu,\nu)$. Apply this to
$S^\gamma$ and $C^\gamma$ and to the laws of $B_{r+1}(H,V)$ and $B_{r+1}(\tree_K,o)$, whose
distance is at most $\eps_1(q,K,r+1)$ by the last statement of Theorem~\ref{thm:localconv}.
\end{proof}

No continuity of the rule is used: the conditional expectations $g_r$ are only measurable,
and this is the reason for proving Theorem~\ref{thm:localconv} in total variation. Agents
whose ball is not a hypertree receive slack $1$; by Theorem~\ref{thm:localconv} their
expected fraction is at most $\eps_1$, and it is included in (d).

\subsection{The split}\label{sec:lp-split}

\begin{lemma}\label{lem:lp-split}
Let $\eta\in(0,1)$ and let $u=(u_p)$ and $w=(w_p)$ be two independent families of
independent uniform variables on $[0,1]$, indexed by the $3q^2$ pairs. Put
\begin{gather*}
\defi'(p)=1-(1-u_p)^{1/(1-\eta)},\qquad\defi''(p)=1-(1-w_p)^{1/\eta},\\
\defi(p)=\min\bigl(\defi'(p),\defi''(p)\bigr).
\end{gather*}
Then the $\defi(p)$ are independent and uniform on $[0,1]$, and for every pair $p$
\[
\defi(p)\le\frac{u_p}{1-\eta},\qquad\defi(p)\le\frac{w_p}{\eta}.
\]
Consequently, writing $\Defi_u$, $\Defi_w$, $\Defi_{\defi}$ for the team deficits computed
from $u$, $w$, $\defi$, we have $\Defi_{\defi}(e)\le\Defi_u(e)/(1-\eta)$ and
$\Defi_{\defi}(e)\le\Defi_w(e)/\eta$ for every team $e$.
\end{lemma}

\begin{proof}
For $t\in[0,1]$, $\P(\defi'(p)>t)=\P(1-u_p<(1-t)^{1-\eta})=(1-t)^{1-\eta}$ and likewise
$\P(\defi''(p)>t)=(1-t)^{\eta}$, so $\P(\defi(p)>t)=1-t$ by independence. By Bernoulli's
inequality $(1-y)^a\ge1-ay$ for $a\ge1$ and $y\in[0,1]$, so $\defi'(p)\le u_p/(1-\eta)$ and
$\defi''(p)\le w_p/\eta$.
\end{proof}

So the pair deficits $\defi$ form a market in the sense of Section~\ref{sec:model}, and so
do $u$ and $w$, which are independent of each other. The main solution will be built on the
uniform market $u$, not on $\defi'$, so local convergence is needed for uniform markets
only; and the comparison of costs is made team by team for a fixed solution, so no
monotonicity of optimal values is needed.

\subsection{The patch}\label{sec:lp-patch}

For a family $\defi$ of pair deficits write $\PhiLP(\defi)$ and $\Phiq(\defi)$ for the
fractional and the integer deficit of the market with these deficits. Recall that
$\E\Phi_n\le C_0\,n^{1/3}$ for every $n\ge1$, by \eqref{eq:variance-C0}, which rests on
\cite[Theorem~2(b)]{FPT20}; we put $\Phi_0=0$.

\begin{lemma}\label{lem:lp-patch}
Let $u$, $w$, $\eta$ and $\defi$ be as in Lemma~\ref{lem:lp-split}. Let $Z$ be a random
element such that the pair $(u,Z)$ is independent of $w$, and let $x=(x_e)$ be a vector of
weights on the teams that is a measurable function of $(u,Z)$, with $x\ge0$ and all loads at
most $1$. Put $\Sigma=q-\sum_ex_e$. Then
\begin{equation}\label{eq:lp-patch}
\E\,\PhiLP(\defi)\ \le\ \frac1{1-\eta}\,\E\sum_e\Defi_u(e)\,x_e+\frac{C_0}{\eta}\,
\bigl(\E\,\Sigma\bigr)^{1/3}.
\end{equation}
If $x$ takes values in $\{0,1\}$, the same bound holds for $\E\,\Phiq(\defi)$.
\end{lemma}

\begin{proof}
Let $\sigma_v=1-\sum_{e\ni v}x_e\in[0,1]$ be the slack of agent $v$. As in
Lemma~\ref{lem:lp-market}(b), the slacks of each sector sum to $\Sigma$.

\emph{Rounding.} Let $U$ be uniform on $[0,1)$ and independent of $(u,Z,w)$. In sector
$i\in\{1,2,3\}$ list the agents as $v_1,\dots,v_q$, put $c_0=0$ and
$c_k=\sigma_{v_1}+\dots+\sigma_{v_k}$, and let
\[
R_i=\bigl\{v_k:\ [c_{k-1},c_k)\text{ contains a point of }U+\{0,1,2,\dots\}\bigr\}.
\]
The interval $[c_{k-1},c_k)$ has length $\sigma_{v_k}\le1$, so it contains at most one such
point, and it contains one with conditional probability $\sigma_{v_k}$ given $(u,Z,w)$. The
intervals partition $[0,\Sigma)$, so
\[
|R_i|=N:=\#\{n\ge0:\ U+n<\Sigma\}
\]
for each of the three sectors. Thus the three sets have the same size
$N\in\{\lfloor\Sigma\rfloor,\lceil\Sigma\rceil\}$, and $\E[N\mid u,Z,w]=\Sigma$.

\emph{The patch.} The triple $R=(R_1,R_2,R_3)$ is a function of $(u,Z,U)$, and $(u,Z,U)$ is
independent of $w$. So, conditionally on $(u,Z,U)$, the variables $w_p$ for the $3N^2$ pairs
inside $R$ are independent and uniform: they form a market with $N$ agents in each sector.
Let $M_R$ be its optimal assignment, a set of $N$ disjoint teams covering exactly the agents
of $R$ (empty if $N=0$). Then
\begin{equation}\label{eq:lp-patchcost}
\E\Bigl[\sum_{e\in M_R}\Defi_w(e)\Bigm|u,Z,U\Bigr]=\E\Phi_n\Big|_{n=N}\le C_0\,N^{1/3}.
\end{equation}
Put $z_e=\P(e\in M_R\mid u,Z,w)$, the average over $U$. Then $z\ge0$, and since an agent lies
in a team of $M_R$ exactly when it belongs to $R$, and then in one team only,
\[
\sum_{e\ni v}z_e=\P(v\in R\mid u,Z,w)=\sigma_v\qquad\text{for every agent }v .
\]
Hence $x+z\ge0$ has all loads exactly $1$: it is a fractional assignment in the sense of
\eqref{eq:model-load}.

\emph{The cost.} By \eqref{eq:model-LP} for the market $\defi$ and by
Lemma~\ref{lem:lp-split}, pointwise,
\[
\PhiLP(\defi)\le\sum_e\Defi_{\defi}(e)\,(x_e+z_e)
\le\frac1{1-\eta}\sum_e\Defi_u(e)\,x_e+\frac1\eta\sum_e\Defi_w(e)\,z_e .
\]
Since $\Defi_w(e)$ is a function of $w$, $\E\sum_e\Defi_w(e)z_e=\E\sum_{e\in M_R}\Defi_w(e)$,
and by \eqref{eq:lp-patchcost} and Jensen's inequality for the concave function
$n\mapsto n^{1/3}$,
\[
\E\sum_{e\in M_R}\Defi_w(e)\le C_0\,\E\,N^{1/3}\le C_0\,(\E N)^{1/3}=C_0\,(\E\,\Sigma)^{1/3}.
\]
This proves \eqref{eq:lp-patch}.

\emph{Integral $x$.} If $x$ takes values in $\{0,1\}$, then $x$ is the indicator of a set of
disjoint teams, $\sigma_v\in\{0,1\}$, $\Sigma$ is an integer, and
$R_i=\{v\text{ in sector }i:\ \sigma_v=1\}$ whatever $U$ is. So $z$ is the indicator of
$M_R$ and $x+z$ is the indicator of a perfect assignment of the market $\defi$. Hence
$\Phiq(\defi)\le\sum_e\Defi_{\defi}(e)(x_e+z_e)$, and the rest of the proof is unchanged.
\end{proof}

The bound \eqref{eq:lp-patch} is concave in $\Sigma$, so only $\E\,\Sigma$ enters. This is
why the convergence of expectations in Lemma~\ref{lem:lp-market}(d) suffices and no
concentration of the slack is needed. The fractional assignment $x+z$ depends on $(u,Z,w)$
and not on $\defi$ alone; this is harmless, since $\PhiLP(\defi)$ is bounded pointwise. The
hypothesis on $Z$ cannot be weakened to independence of $Z$ and $w$: if $Z=(u+w)\bmod1$
coordinatewise, then $Z$ is independent of $w$, but $(u,Z)$ determines $w$. In the proof of
Theorem~\ref{thm:lp} there is no $Z$.

\subsection{Proof of Theorem~\ref{thm:lp}}\label{sec:lp-proof}

\begin{proof}[Proof of Theorem~\ref{thm:lp}]
Let $c_r$, $\eps_r$ and the rules $g_r$ be as in Lemma~\ref{lem:lp-condrule}; they depend
only on the law of $(\tree,o,\xi,m)$, $c_r\le c(m)$ and $\eps_r\to0$. Fix $r\ge1$,
$\eta\in(0,1)$ and $q\ge r^{3/2}$, so that $K=r\le q^{2/3}$. Let $u$, $w$ and $\defi$ be as
in Lemma~\ref{lem:lp-split}. Since $\defi$ is a market, $\E\PhiLP=\E\,\PhiLP(\defi)$. Apply
the rule $g_r$ to the near-tight hypergraph $H=H_q(r)$ of the market $u$, and let
$x=x^{g_r}$. By Lemma~\ref{lem:lp-market}(a) and~(b), $x$ is a function of $u$ that
satisfies the hypotheses of Lemma~\ref{lem:lp-patch}, with trivial $Z$. By
\eqref{eq:lp-patch} and Lemma~\ref{lem:lp-market}(c),
\[
\frac{\E\PhiLP}{q^{1/3}}\le\frac{\E\,C^{g_r}(B_{r+1}(H,V))}{1-\eta}
+\frac{C_0}{\eta}\Bigl(\E\,S^{g_r}(B_{r+1}(H,V))\Bigr)^{1/3},
\]
and by Lemma~\ref{lem:lp-market}(d),
\begin{align*}
\E\,C^{g_r}(B_{r+1}(H,V))&\le c_r+r\,\eps_1(q,r,r+1),\\
\E\,S^{g_r}(B_{r+1}(H,V))&\le\eps_r+\eps_1(q,r,r+1).
\end{align*}
This is \eqref{eq:lp-explicit}.

\emph{Order of limits.} Let $q\to\infty$ with $r$ and $\eta$ fixed. By
Theorem~\ref{thm:localconv}, $\eps_1(q,r,r+1)\to0$, so
\[
\limsup_{q\to\infty}\frac{\E\PhiLP}{q^{1/3}}\le\frac{c_r}{1-\eta}+\frac{C_0}{\eta}\,
\eps_r^{1/3}\le\frac{c(m)}{1-\eta}+\frac{C_0}{\eta}\,\eps_r^{1/3}.
\]
The left-hand side does not depend on $r$ or $\eta$. Let $r\to\infty$, and then
$\eta\to0$.
\end{proof}

\begin{remark}\label{rem:lp-scope}
The theorem is a statement about expectations, and it gives no value of $c(m)$: an invariant
perfect matching of small cost is supplied by Theorems~\ref{thm:process}
and~\ref{thm:existence}. The proof uses $m$ only through the conditional expectations
$h_j$. Unbounded costs never appear in it, because the rule $g_r$ ignores the teams of cost
above $r$. The bound \eqref{eq:lp-explicit} is explicit in $q$ apart from the quantities
$c_r$ and $\eps_r$ of the tree, but it is of no use at any size that can be simulated; see
Remark~\ref{rem:localconv-size}.
\end{remark}

\begin{remark}[Floating point]\label{rem:lp-check}
The rounding and the patch of Lemma~\ref{lem:lp-patch} were checked numerically on small
instances by the script \texttt{patch\_check.py} in the folder
\texttt{lemmaC/round3/lplocal/} ($q=5$, $\eta=0.25$, six seeds, a random packing $x$, exact integration over $U$, and $M_R$ found by exhaustive
search). In all runs the three sector slack sums agreed to $10^{-12}$, the three sets
$R_i$ had the same size in $\{\lfloor\Sigma\rfloor,\lceil\Sigma\rceil\}$, the loads of $x+z$
differed from $1$ by at most $1.1\cdot10^{-15}$, and the chain of inequalities for the cost
held. This is evidence only and is not used in the proof.
\end{remark}
